% =====================================================================
%  CP Reference Book - C++ Edition
%  Adan Arteaga
%  Self-contained: drop this single file into Overleaf and it compiles.
%  Build: pdflatex -interaction=nonstopmode cp-reference-cpp.tex  (twice)
% =====================================================================
\documentclass[10pt]{article}

% ---------------------------------------------------------------------
% Language and encoding
% ---------------------------------------------------------------------
\usepackage[english]{babel}
\usepackage[utf8]{inputenc}
\usepackage[T1]{fontenc}
\usepackage{lmodern}

% ---------------------------------------------------------------------
% Page geometry: landscape, tight margins, two columns
% ---------------------------------------------------------------------
\usepackage[landscape,margin=1cm,top=1.4cm,bottom=1.2cm,headsep=0.35cm]{geometry}
\usepackage{multicol}
\setlength{\columnsep}{0.75cm}
\setlength{\columnseprule}{0.4pt}
\setlength{\parindent}{0pt}
\setlength{\parskip}{2.5pt}
\raggedcolumns

% ---------------------------------------------------------------------
% Core packages
% ---------------------------------------------------------------------
\usepackage{amsmath}
\usepackage{amssymb}
\usepackage{xcolor}
\usepackage{enumitem}
\usepackage{booktabs}
\usepackage{array}
\usepackage{listings}
\usepackage{fancyhdr}
\usepackage[nobottomtitles*]{titlesec}
\usepackage{hyperref}

% ---------------------------------------------------------------------
% Palette
% ---------------------------------------------------------------------
\definecolor{accent}{HTML}{18456B}   % section bars, headings
\definecolor{accent2}{HTML}{0E7C66}  % code rule, complexity badges
\definecolor{codebg}{HTML}{F5F7F9}   % code background
\definecolor{hairline}{HTML}{C3CFD9} % rules and separators
\definecolor{kw}{HTML}{0B5FA5}       % keywords
\definecolor{cmt}{HTML}{6A737D}      % comments and muted text
\definecolor{strc}{HTML}{1B7A3D}     % strings
\definecolor{warnc}{HTML}{B00020}    % warnings
\definecolor{tipc}{HTML}{9A6100}     % tips

\renewcommand{\columnseprulecolor}{\color{hairline}}

\hypersetup{
    colorlinks=true,
    linkcolor=accent,
    urlcolor=accent2,
    bookmarksnumbered=true,
    bookmarksopen=true,
    bookmarksdepth=2,
    pdftitle={CP Reference Book - C++ Edition},
    pdfauthor={Adan Arteaga}
}

% ---------------------------------------------------------------------
% Headings
% ---------------------------------------------------------------------
\setcounter{secnumdepth}{2}
\setcounter{tocdepth}{1}          % printed TOC: sections only, stays on one page
\renewcommand{\bottomtitlespace}{0.09\textheight}

% titlesec passes the title to the LAST command of the before-code,
% so both headings are wrapped in a helper that takes it as #1.
\newcommand{\sectionbar}[1]{%
  \hspace*{-\fboxsep}\colorbox{accent}{%
    \parbox{\dimexpr\linewidth-2\fboxsep\relax}{%
      \vspace{2.5pt}\color{white}\thesection\hspace{0.8em}#1\vspace{2.5pt}}}}
% Starred sections (the TOC heading) must not print a number.
\newcommand{\sectionbarplain}[1]{%
  \hspace*{-\fboxsep}\colorbox{accent}{%
    \parbox{\dimexpr\linewidth-2\fboxsep\relax}{%
      \vspace{2.5pt}\color{white}#1\vspace{2.5pt}}}}
% Thin rule under the title, drawn inside the heading so it cannot
% collide with the text the way a titlesec after-code rule does.
\newcommand{\subhead}[1]{%
  \raisebox{0.35ex}{\color{accent2}\rule{2.1mm}{2.1mm}}\hspace{0.55em}#1%
  \par\nobreak\vspace{1.5pt}{\color{hairline}\hrule height 0.4pt}\nobreak}

\titleformat{\section}[block]
  {\normalfont\large\bfseries}{}{0pt}{\sectionbar}
\titleformat{name=\section,numberless}[block]
  {\normalfont\large\bfseries}{}{0pt}{\sectionbarplain}
\titlespacing*{\section}{0pt}{16pt}{9pt}

\titleformat{\subsection}[block]
  {\normalfont\normalsize\bfseries\color{accent}}{}{0pt}{\subhead}
\titlespacing*{\subsection}{0pt}{8pt}{-7pt}

\setlist[itemize]{leftmargin=1.1em, itemsep=1pt, topsep=2pt, parsep=0pt}
\setlist[enumerate]{leftmargin=1.4em, itemsep=1pt, topsep=2pt, parsep=0pt}

% ---------------------------------------------------------------------
% Running head / foot
% ---------------------------------------------------------------------
\pagestyle{fancy}
\fancyhf{}
\fancyhead[L]{\footnotesize\color{cmt}\nouppercase{\leftmark}}
\fancyhead[R]{\footnotesize\color{cmt}CP Reference Book \textbullet\ C++ Edition}
\fancyfoot[C]{\footnotesize\color{cmt}\thepage}
\renewcommand{\headrulewidth}{0.4pt}
\renewcommand{\headrule}{\color{hairline}\hrule height \headrulewidth}

% ---------------------------------------------------------------------
% Code style
% ---------------------------------------------------------------------
\lstdefinestyle{cpp}{
    language=C++,
    basicstyle=\ttfamily\scriptsize,
    backgroundcolor=\color{codebg},
    keywordstyle=\bfseries\color{kw},
    commentstyle=\itshape\color{cmt},
    stringstyle=\color{strc},
    frame=leftline,
    framerule=1.3pt,
    rulecolor=\color{accent2},
    framexleftmargin=4pt,
    framexrightmargin=2pt,
    xleftmargin=5pt,
    breaklines=true,
    breakatwhitespace=false,
    postbreak=\mbox{\textcolor{cmt}{$\hookrightarrow$}\space},
    columns=fullflexible,
    keepspaces=true,
    showstringspaces=false,
    tabsize=2,
    aboveskip=5pt,
    belowskip=5pt,
    upquote=true,
    morekeywords={ll,pii,vi,int64_t,uint64_t,size_t,constexpr,nullptr,auto}
}
\lstset{style=cpp}

% ---------------------------------------------------------------------
% Semantic macros
% ---------------------------------------------------------------------
\newcommand{\code}[1]{\texttt{\small #1}}
% \sub{Title}{complexity without \mathcal{O}} -> heading + complexity badge
\newcommand{\sub}[2]{\subsection[#1]{#1\hspace{0.6em}%
  \textnormal{\color{accent2}\small$\mathcal{O}(#2)$}}}
\newcommand{\uses}[1]{{\color{accent2}\footnotesize\textbf{Use:}}~{\footnotesize #1}\par}
\newcommand{\tip}[1]{{\color{tipc}\footnotesize\textbf{Tip:}}~{\footnotesize #1}\par}
\newcommand{\warn}[1]{{\color{warnc}\footnotesize\textbf{Careful:}}~{\footnotesize #1}\par}
\newcommand{\note}[1]{{\footnotesize\color{cmt}#1}\par}

\begin{document}
\pagenumbering{gobble}

% =====================================================================
% Cover
% =====================================================================
\begin{titlepage}
    \centering
    \vspace*{1.6cm}

    {\color{accent}\rule{\linewidth}{2.5pt}}
    \vspace{0.7cm}

    {\Huge\bfseries\color{accent} CP Reference Book\par}
    \vspace{0.5cm}
    {\LARGE C++ Edition\par}

    \vspace{0.5cm}
    {\color{accent}\rule{\linewidth}{2.5pt}}

    \vspace{1.4cm}
    {\Large Adán Arteaga\par}

    \vspace{1.4cm}
    {\large A contest-ready reference for competitive programming\\
    and technical interviews\par}

    \vspace{0.9cm}
    {\normalsize\color{cmt}
    Every snippet compiles and runs as written --- no pseudocode,\\
    no placeholders. Complexities are stated where the choice matters.\par}

    \vfill
    {\footnotesize\color{cmt}Landscape layout, two columns, built for printing.\\
    Clickable table of contents and full PDF bookmarks for on-screen use.\par}
\end{titlepage}

\pagenumbering{roman}
\thispagestyle{plain}
\tableofcontents
\newpage
\pagenumbering{arabic}

% =====================================================================
\section{Contest Setup and Complexity}
% =====================================================================
\begin{multicols}{2}

\subsection{Full Template}
\begin{lstlisting}
#include <bits/stdc++.h>
using namespace std;

using ll  = long long;
using ull = unsigned long long;
using pii = pair<int,int>;
using pll = pair<ll,ll>;

#define all(x)  (x).begin(), (x).end()
#define rall(x) (x).rbegin(), (x).rend()
#define sz(x)   (int)(x).size()

const ll INF  = 4e18;
const int IINF = 1e9;
const ll MOD  = 1e9 + 7;

void solve() {
    // ...
}

int main() {
    ios::sync_with_stdio(false);
    cin.tie(nullptr);

    int t = 1;
    cin >> t;              // remove if single test case
    while (t--) solve();
    return 0;
}
\end{lstlisting}
\warn{\code{bits/stdc++.h} is GCC-only. On MSVC or Clang without it,
include the headers you need explicitly.}

\subsection{Fast I/O}
\begin{lstlisting}
ios::sync_with_stdio(false);
cin.tie(nullptr);
\end{lstlisting}
\warn{After this you cannot mix \code{cin} with \code{scanf}, or
\code{cout} with \code{printf}. Pick one family and stay in it.}

\tip{\code{'\textbackslash n'} is faster than \code{endl}: \code{endl}
flushes the buffer every single time. Only flush deliberately (interactive
problems).}

\subsection{Decimal Precision}
\begin{lstlisting}
cout << fixed << setprecision(10) << ans << '\n';
\end{lstlisting}
\tip{Compare doubles with an epsilon, never with \code{==}:
\code{fabs(a - b) < 1e-9}. When the answer is an integer computed in
floating point, round with \code{llround(x)} instead of casting.}

\subsection{Data Type Limits}
\begin{center}\footnotesize
\begin{tabular}{@{}lll@{}}
\toprule
\textbf{Type} & \textbf{Max} & \textbf{Precision} \\
\midrule
\code{int}          & $\approx 2\cdot 10^{9}$   & --- \\
\code{unsigned}     & $\approx 4\cdot 10^{9}$   & --- \\
\code{long long}    & $\approx 9\cdot 10^{18}$  & --- \\
\code{unsigned ll}  & $\approx 1.8\cdot 10^{19}$& --- \\
\code{\_\_int128}   & $\approx 1.7\cdot 10^{38}$& GCC only \\
\code{float}        & ---                       & 7 digits \\
\code{double}       & ---                       & 15--17 digits \\
\code{long double}  & ---                       & 18--19 digits \\
\bottomrule
\end{tabular}
\end{center}
\warn{The single most common bug in contests: \code{int * int} overflows
\emph{before} it is assigned to a \code{long long}. Cast one operand
first: \code{(ll)a * b}.}

\subsection{Printing \_\_int128}
\begin{lstlisting}
void print128(__int128 x) {
    if (x < 0) { putchar('-'); x = -x; }
    if (x > 9) print128(x / 10);
    putchar((char)('0' + (int)(x % 10)));
}
\end{lstlisting}
\uses{exact products of two \code{long long} values, and safe modular
multiplication when \code{MOD} exceeds $2^{31}$.}

\subsection{Complexity Classes}
\begin{itemize}
    \item $O(1)$ --- array index, hash lookup on average.
    \item $O(\log n)$ --- binary search, \code{set}, heap push/pop.
    \item $O(\sqrt n)$ --- primality, divisor enumeration.
    \item $O(n)$ --- a single pass.
    \item $O(n \log n)$ --- sorting, most greedy solutions.
    \item $O(n^2)$ --- nested loops, many DP tables.
    \item $O(2^n)$ --- subsets, bitmask DP.
    \item $O(n!)$ --- permutations.
\end{itemize}

\subsection{Constraint to Complexity}
Assume roughly $10^8$ simple operations per second.
\begin{center}\footnotesize
\begin{tabular}{@{}ll@{}}
\toprule
\textbf{Constraint} & \textbf{Target} \\
\midrule
$n \le 10$        & $O(n!)$ \\
$n \le 20$        & $O(2^n \cdot n)$ \\
$n \le 100$       & $O(n^4)$ \\
$n \le 500$       & $O(n^3)$ \\
$n \le 5000$      & $O(n^2)$ \\
$n \le 10^5$      & $O(n \log n)$ \\
$n \le 10^6$      & $O(n)$ or $O(n \log n)$ \\
$n \le 10^{18}$   & $O(\log n)$ or $O(\sqrt n)$ \\
\bottomrule
\end{tabular}
\end{center}
\tip{Read the constraints \emph{before} designing the solution. They
usually name the intended complexity, which names the algorithm.}

\subsection{STL Complexity}
\textbf{vector}
\begin{itemize}
    \item \code{push\_back}, \code{pop\_back}, \code{[]} $\rightarrow O(1)$
          amortized.
    \item \code{insert}, \code{erase} in the middle $\rightarrow O(n)$.
\end{itemize}

\textbf{set / map / multiset} (balanced BST)
\begin{itemize}
    \item \code{insert}, \code{erase}, \code{find},
          \code{lower\_bound} $\rightarrow O(\log n)$.
\end{itemize}

\textbf{unordered\_map / unordered\_set} (hash table)
\begin{itemize}
    \item \code{insert}, \code{find} $\rightarrow O(1)$ average,
          $O(n)$ worst case.
\end{itemize}

\textbf{priority\_queue} (binary heap)
\begin{itemize}
    \item \code{push}, \code{pop} $\rightarrow O(\log n)$;
          \code{top} $\rightarrow O(1)$.
\end{itemize}

\textbf{deque}
\begin{itemize}
    \item \code{push\_front}, \code{push\_back}, \code{pop} on both
          ends $\rightarrow O(1)$.
\end{itemize}

\subsection{Memory Budget}
\begin{lstlisting}
// 4 bytes each: 64M ints ~= 256 MB
int a[64 * 1024 * 1024];
\end{lstlisting}
\tip{Global arrays are zero-initialised and live in static storage;
huge local arrays overflow the stack. Declare big arrays at global scope.}

\subsection{Compilation and Debugging}
\begin{lstlisting}
g++ -std=c++17 -O2 -Wall -Wextra sol.cpp -o sol

// Catches out-of-bounds and UB at runtime:
g++ -std=c++17 -g -fsanitize=address,undefined sol.cpp
\end{lstlisting}

\begin{lstlisting}
// Debug print that vanishes on the judge
#ifdef LOCAL
  #define dbg(x) cerr << #x << " = " << (x) << '\n'
#else
  #define dbg(x)
#endif
\end{lstlisting}
\tip{Compile locally with \code{-DLOCAL}. \code{cerr} is not judged, so
leftover debug output on \code{cerr} will not cause a wrong answer ---
but it does cost time.}

\end{multicols}

% =====================================================================
\section{C++ STL Toolkit}
% =====================================================================
\begin{multicols}{2}

\subsection{Vector Essentials}
\begin{lstlisting}
vector<int> v(n);              // n zeros
vector<int> v(n, -1);          // n copies of -1
vector<vector<int>> g(n, vector<int>(m, 0));

sort(all(v));
sort(rall(v));                 // descending
reverse(all(v));

// Remove duplicates (must be sorted first)
sort(all(v));
v.erase(unique(all(v)), v.end());

int mn = *min_element(all(v));
int mx = *max_element(all(v));
ll  s  = accumulate(all(v), 0LL);   // 0LL avoids int overflow
int c  = count(all(v), x);
\end{lstlisting}
\warn{\code{accumulate(all(v), 0)} sums into an \code{int}. Always pass
\code{0LL} when the total can exceed $2\cdot 10^9$.}

\subsection{Sorting with Custom Order}
\begin{lstlisting}
// By second value, then first descending
sort(all(v), [](const pii& a, const pii& b) {
    if (a.second != b.second) return a.second < b.second;
    return a.first > b.first;
});

// Sort indices instead of moving heavy objects
vector<int> idx(n);
iota(all(idx), 0);
sort(all(idx), [&](int i, int j) { return a[i] < a[j]; });
\end{lstlisting}
\warn{A comparator must be a \emph{strict} weak ordering: it must return
\code{false} for equal elements. Using \code{<=} is undefined behaviour
and can segfault.}

\sub{nth\_element}{n}
\begin{lstlisting}
// Places the k-th smallest at position k,
// partitions around it. No full sort.
nth_element(v.begin(), v.begin() + k, v.end());
int kth = v[k];
\end{lstlisting}

\subsection{Pairs, Tuples and Bindings}
\begin{lstlisting}
pair<int,int> p = {1, 2};
tuple<int,int,int> t = {1, 2, 3};

auto [x, y] = p;               // C++17
auto [a, b, c] = t;

for (auto& [key, val] : mp) {  // iterate a map cleanly
    // ...
}
\end{lstlisting}

\subsection{set / map Operations}
\begin{lstlisting}
set<int> s;
s.insert(x);
s.erase(x);                    // by value
s.erase(s.begin());            // by iterator

// First element >= x  /  first element > x
auto it = s.lower_bound(x);
auto jt = s.upper_bound(x);
if (it != s.end()) { /* *it exists */ }

// Largest element <= x
auto kt = s.upper_bound(x);
if (kt != s.begin()) { --kt; /* *kt <= x */ }
\end{lstlisting}
\warn{Use the \emph{member} \code{s.lower\_bound(x)}, not
\code{std::lower\_bound(all(s), x)}. The free function on a \code{set}
degrades to $O(n)$ because the iterators are not random access.}

\begin{lstlisting}
map<int,int> mp;
mp[k]++;                       // inserts 0 then increments
if (mp.count(k)) { /* ... */ }
mp.erase(k);

// Does NOT insert when missing
auto it = mp.find(k);
if (it != mp.end()) use(it->second);
\end{lstlisting}
\warn{\code{mp[k]} default-constructs a missing key. Inside a read-only
loop that silently grows the map; use \code{find} or \code{count}.}

\subsection{multiset: Erase One vs All}
\begin{lstlisting}
multiset<int> ms;
ms.erase(x);                   // erases EVERY copy of x
ms.erase(ms.find(x));          // erases exactly one
\end{lstlisting}

\subsection{priority\_queue}
\begin{lstlisting}
priority_queue<int> maxHeap;   // largest on top (default)

priority_queue<int, vector<int>, greater<int>> minHeap;

// Min-heap of pairs, the common Dijkstra shape
priority_queue<pll, vector<pll>, greater<pll>> pq;
pq.push({dist, node});
\end{lstlisting}
\tip{No \code{greater} handy? Push negated values into a max-heap and
negate on the way out.}

\subsection{deque}
\begin{lstlisting}
deque<int> dq;
dq.push_back(x);  dq.push_front(x);
dq.pop_back();    dq.pop_front();
int f = dq.front(), b = dq.back();
\end{lstlisting}
\uses{sliding window extrema and 0--1 BFS.}

\subsection{bitset}
\begin{lstlisting}
bitset<1000> b;
b[i] = 1;
b.set();  b.reset();  b.flip();
int ones = b.count();
bool any = b.any();

// Set operations run 64 bits at a time
bitset<1000> c = a & b;
c = a | b;   c = a ^ b;   c = a << 3;
\end{lstlisting}
\uses{a constant-factor $1/64$ speedup. It turns an $O(n^2)$ DP or
reachability computation into $O(n^2/64)$, which is often enough to
pass.}

\subsection{next\_permutation}
\begin{lstlisting}
sort(all(v));                  // must start sorted
do {
    // process this permutation
} while (next_permutation(all(v)));
\end{lstlisting}

\subsection{Anti-Hash Custom Hash}
\begin{lstlisting}
struct chash {
    static uint64_t splitmix64(uint64_t x) {
        x += 0x9e3779b97f4a7c15ULL;
        x = (x ^ (x >> 30)) * 0xbf58476d1ce4e5b9ULL;
        x = (x ^ (x >> 27)) * 0x94d049bb133111ebULL;
        return x ^ (x >> 31);
    }
    size_t operator()(uint64_t x) const {
        static const uint64_t SEED =
            chrono::steady_clock::now()
                .time_since_epoch().count();
        return splitmix64(x + SEED);
    }
};

unordered_map<ll, int, chash> mp;
mp.reserve(1 << 20);
mp.max_load_factor(0.25);
\end{lstlisting}
\warn{The default \code{unordered\_map} hash for integers is the
identity. On Codeforces, anyone can feed you colliding keys and turn
your $O(1)$ lookups into $O(n)$. Use this on any open hacking round, or
just use \code{map}.}

\subsection{Randomness}
\begin{lstlisting}
mt19937_64 rng(chrono::steady_clock::now()
                   .time_since_epoch().count());

ll r = rng() % n;                     // biased but usually fine
ll u = uniform_int_distribution<ll>(1, n)(rng);
shuffle(all(v), rng);
\end{lstlisting}
\warn{Never seed with a constant, and never use \code{rand()} --- on
Windows it only produces values up to $32767$.}

\end{multicols}

% =====================================================================
\section{Math, Modular Arithmetic and Number Theory}
% =====================================================================
\begin{multicols}{2}

\sub{GCD and LCM}{\log \min(a,b)}
\begin{lstlisting}
// GCC built-in, needs <algorithm>
int g = __gcd(a, b);

ll gcd_it(ll a, ll b) {
    while (b) { a %= b; swap(a, b); }
    return a < 0 ? -a : a;
}

ll lcm_safe(ll a, ll b) {
    if (a == 0 || b == 0) return 0;
    return (a / __gcd(a, b)) * b;   // divide FIRST
}
\end{lstlisting}
\textbf{Properties:}
\begin{itemize}
    \item $\gcd(a,b)\cdot\mathrm{lcm}(a,b) = |a\cdot b|$.
    \item $\gcd(a,b) = \gcd(b, a \bmod b)$.
    \item $\gcd(a,b,c) = \gcd(a, \gcd(b,c))$.
\end{itemize}
\warn{In \code{lcm}, dividing before multiplying is what keeps the
product from overflowing.}

\sub{Extended Euclid}{\log \min(a,b)}
Solves $ax + by = \gcd(a,b)$.
\begin{lstlisting}
ll extgcd(ll a, ll b, ll& x, ll& y) {
    if (b == 0) { x = 1; y = 0; return a; }
    ll x1, y1;
    ll g = extgcd(b, a % b, x1, y1);
    x = y1;
    y = x1 - (a / b) * y1;
    return g;
}
\end{lstlisting}
\uses{modular inverse when \code{MOD} is \emph{not} prime, and linear
Diophantine equations.}

\subsection{Linear Diophantine}
$ax + by = c$ has a solution iff $\gcd(a,b) \mid c$.
\begin{lstlisting}
bool solve_dio(ll a, ll b, ll c, ll& x, ll& y) {
    ll g = extgcd(a, b, x, y);
    if (c % g != 0) return false;
    x *= c / g;
    y *= c / g;
    return true;
}
// General family:
// x + k*(b/g),  y - k*(a/g)   for any integer k
\end{lstlisting}

\subsection{Modulo Properties}
\[(a+b) \bmod m = ((a \bmod m)+(b \bmod m)) \bmod m\]
\[(a-b) \bmod m = ((a \bmod m)-(b \bmod m)+m) \bmod m\]
\[(a \cdot b) \bmod m = ((a \bmod m)(b \bmod m)) \bmod m\]
\begin{lstlisting}
// C++ % keeps the sign of the dividend
int r = (a % MOD + MOD) % MOD;
\end{lstlisting}
\warn{Division does \emph{not} distribute over the modulus. Multiply by
the modular inverse instead.}

\sub{Binary Exponentiation}{\log b}
\begin{lstlisting}
ll binpow(ll a, ll b, ll mod) {
    ll res = 1;
    a %= mod;
    if (a < 0) a += mod;
    while (b > 0) {
        if (b & 1) res = res * a % mod;
        a = a * a % mod;
        b >>= 1;
    }
    return res;
}
\end{lstlisting}

\subsection{Safe Multiplication mod m}
\begin{lstlisting}
// Needed when mod > 2^31 (a*b overflows ll)
ll mulmod(ll a, ll b, ll mod) {
    return (__int128)a * b % mod;
}
\end{lstlisting}

\sub{Modular Inverse}{\log m}
\begin{lstlisting}
// MOD prime, a not divisible by MOD (Fermat)
ll inv(ll a, ll mod) { return binpow(a, mod - 2, mod); }

// Any mod, requires gcd(a, mod) == 1
ll inv_gen(ll a, ll mod) {
    ll x, y;
    ll g = extgcd(a, mod, x, y);
    if (g != 1) return -1;         // no inverse
    return (x % mod + mod) % mod;
}
\end{lstlisting}

\sub{All Inverses 1..n}{n}
\begin{lstlisting}
vector<ll> inv(n + 1);
inv[1] = 1;
for (int i = 2; i <= n; i++)
    inv[i] = (MOD - (MOD / i) * inv[MOD % i] % MOD) % MOD;
\end{lstlisting}
\uses{one linear pass instead of $n$ separate $O(\log m)$ exponentiations.}

\sub{Factorials and nCr}{n \text{ precompute},\; 1 \text{ per query}}
\begin{lstlisting}
const int MAXN = 1e6 + 5;
ll fact[MAXN], invfact[MAXN];

void init_fact() {
    fact[0] = 1;
    for (int i = 1; i < MAXN; i++)
        fact[i] = fact[i - 1] * i % MOD;

    invfact[MAXN - 1] = binpow(fact[MAXN - 1], MOD - 2, MOD);
    for (int i = MAXN - 1; i > 0; i--)
        invfact[i - 1] = invfact[i] * i % MOD;
}

ll nCr(ll n, ll r) {
    if (r < 0 || r > n) return 0;
    return fact[n] * invfact[r] % MOD * invfact[n - r] % MOD;
}

ll nPr(ll n, ll r) {
    if (r < 0 || r > n) return 0;
    return fact[n] * invfact[n - r] % MOD;
}
\end{lstlisting}
\[C(n,r) = \frac{n!}{r!\,(n-r)!} \qquad P(n,r) = \frac{n!}{(n-r)!}\]
\tip{The backward loop computes every inverse factorial with a single
exponentiation. Calling \code{binpow} inside \code{nCr} is the slow way
and shows up as TLE with many queries.}

\subsection{Combinatorial Identities}
\begin{itemize}
    \item $C(n,r) = C(n-1,r-1) + C(n-1,r)$ --- Pascal.
    \item $\sum_{r=0}^{n} C(n,r) = 2^n$.
    \item Stars and bars: $x_1+\dots+x_k=n$, $x_i \ge 0$
          $\rightarrow C(n+k-1,\,k-1)$.
    \item Catalan: $C_n = \frac{1}{n+1}\binom{2n}{n}$ --- balanced
          brackets, binary trees, triangulations.
\end{itemize}

\sub{Sieve of Eratosthenes}{n \log\log n}
\begin{lstlisting}
const int MAXN = 1e6 + 5;
vector<bool> isPrime(MAXN, true);
vector<int> primes;

void sieve(int n) {
    isPrime[0] = isPrime[1] = false;
    for (int i = 2; i <= n; i++) {
        if (isPrime[i]) {
            primes.push_back(i);
            for (ll j = (ll)i * i; j <= n; j += i)
                isPrime[j] = false;
        }
    }
}
\end{lstlisting}
\warn{Start the inner loop at $i^2$ and use \code{ll} for \code{j}:
$i \cdot i$ overflows \code{int} once $i > 46340$.}

\sub{Linear Sieve with SPF}{n}
\begin{lstlisting}
vector<int> spf(MAXN, 0), primes;

void linear_sieve(int n) {
    for (int i = 2; i <= n; i++) {
        if (spf[i] == 0) {
            spf[i] = i;
            primes.push_back(i);
        }
        for (int p : primes) {
            if (p > spf[i] || (ll)i * p > n) break;
            spf[i * p] = p;
        }
    }
}

// O(log x) factorization after the sieve
vector<int> factorize_spf(int x) {
    vector<int> f;
    while (x > 1) { f.push_back(spf[x]); x /= spf[x]; }
    return f;
}
\end{lstlisting}
\uses{when you must factor many numbers. One $O(n)$ precomputation
turns each factorization into $O(\log x)$.}

\sub{Trial Division Factorization}{\sqrt n}
\begin{lstlisting}
vector<pll> factorize(ll x) {
    vector<pll> f;                 // {prime, exponent}
    for (ll p = 2; p * p <= x; p++) {
        if (x % p == 0) {
            int e = 0;
            while (x % p == 0) { x /= p; e++; }
            f.push_back({p, e});
        }
    }
    if (x > 1) f.push_back({x, 1});
    return f;
}
\end{lstlisting}

\sub{All Divisors}{\sqrt n}
\begin{lstlisting}
vector<ll> divisors(ll n) {
    vector<ll> d;
    for (ll i = 1; i * i <= n; i++) {
        if (n % i == 0) {
            d.push_back(i);
            if (i != n / i) d.push_back(n / i);
        }
    }
    sort(all(d));
    return d;
}
\end{lstlisting}
\textbf{From the factorization $n=\prod p_i^{e_i}$:}
\begin{itemize}
    \item Number of divisors: $\prod (e_i + 1)$.
    \item Sum of divisors: $\prod \frac{p_i^{e_i+1}-1}{p_i-1}$.
\end{itemize}

\sub{Divisor Sieve}{n \log n}
\begin{lstlisting}
// divs[i] = list of all divisors of i, for every i <= n
vector<vector<int>> divs(n + 1);
for (int i = 1; i <= n; i++)
    for (int j = i; j <= n; j += i)
        divs[j].push_back(i);
\end{lstlisting}

\subsection{Euler's Totient Function}
$\phi(n)$ counts integers in $[1,n]$ coprime to $n$.
\begin{lstlisting}
// Single value, O(sqrt n)
ll phi(ll n) {
    ll res = n;
    for (ll p = 2; p * p <= n; p++) {
        if (n % p == 0) {
            while (n % p == 0) n /= p;
            res -= res / p;
        }
    }
    if (n > 1) res -= res / n;
    return res;
}

// Sieve for all i <= n, O(n log log n)
vector<int> phi_all(int n) {
    vector<int> f(n + 1);
    iota(f.begin(), f.end(), 0);
    for (int i = 2; i <= n; i++)
        if (f[i] == i)                 // i is prime
            for (int j = i; j <= n; j += i)
                f[j] -= f[j] / i;
    return f;
}
\end{lstlisting}
\uses{Euler's theorem: $a^{\phi(m)} \equiv 1 \pmod m$ when
$\gcd(a,m)=1$. This is how you reduce huge exponents modulo a composite
$m$.}

\sub{Miller-Rabin Primality}{k \log^3 n}
\begin{lstlisting}
bool check_composite(ull n, ull a, ull d, int s) {
    ull x = binpow(a, d, n);
    if (x == 1 || x == n - 1) return false;
    for (int r = 1; r < s; r++) {
        x = (__int128)x * x % n;
        if (x == n - 1) return false;
    }
    return true;
}

bool is_prime(ull n) {
    if (n < 2) return false;
    int s = 0;
    ull d = n - 1;
    while ((d & 1) == 0) { d >>= 1; s++; }

    // Deterministic for every 64-bit n
    for (ull a : {2ULL, 3ULL, 5ULL, 7ULL, 11ULL, 13ULL,
                  17ULL, 19ULL, 23ULL, 29ULL, 31ULL, 37ULL}) {
        if (n == a) return true;
        if (check_composite(n, a, d, s)) return false;
    }
    return true;
}
\end{lstlisting}
\uses{primality of values up to $10^{18}$, where trial division is
hopeless. Requires the \code{\_\_int128} \code{binpow}.}

\sub{Pollard's Rho Factorization}{n^{1/4}}
\begin{lstlisting}
ull pollard(ull n) {
    if (n % 2 == 0) return 2;
    ull x = rng() % n, y = x, c = rng() % n + 1, d = 1;
    while (d == 1) {
        x = ((__int128)x * x + c) % n;
        y = ((__int128)y * y + c) % n;
        y = ((__int128)y * y + c) % n;
        d = __gcd(x > y ? x - y : y - x, n);
    }
    return d == n ? pollard(n) : d;
}

void factor(ull n, vector<ull>& out) {
    if (n == 1) return;
    if (is_prime(n)) { out.push_back(n); return; }
    ull d = pollard(n);
    factor(d, out);
    factor(n / d, out);
}
\end{lstlisting}
\uses{factoring numbers up to $10^{18}$ in milliseconds.}

\sub{Chinese Remainder Theorem}{\log m}
Merges $x \equiv a_1 \pmod{m_1}$ and $x \equiv a_2 \pmod{m_2}$.
\begin{lstlisting}
// Returns {x, lcm} or {-1, -1} when incompatible
pll crt(ll a1, ll m1, ll a2, ll m2) {
    ll p, q;
    ll g = extgcd(m1, m2, p, q);
    if ((a2 - a1) % g != 0) return {-1, -1};

    ll lcm = m1 / g * m2;
    __int128 t = (__int128)(a2 - a1) / g % (m2 / g) * p % (m2 / g);
    ll x = (ll)(((__int128)a1 + (__int128)t * m1) % lcm);
    return {(x % lcm + lcm) % lcm, lcm};
}
\end{lstlisting}

\sub{Matrix Exponentiation}{k^3 \log n}
\begin{lstlisting}
using Mat = vector<vector<ll>>;

Mat mul(const Mat& A, const Mat& B) {
    int n = A.size(), m = B[0].size(), k = B.size();
    Mat C(n, vector<ll>(m, 0));
    for (int i = 0; i < n; i++)
        for (int p = 0; p < k; p++) {
            if (!A[i][p]) continue;
            for (int j = 0; j < m; j++)
                C[i][j] = (C[i][j] + A[i][p] * B[p][j]) % MOD;
        }
    return C;
}

Mat matpow(Mat A, ll e) {
    int n = A.size();
    Mat R(n, vector<ll>(n, 0));
    for (int i = 0; i < n; i++) R[i][i] = 1;
    while (e) {
        if (e & 1) R = mul(R, A);
        A = mul(A, A);
        e >>= 1;
    }
    return R;
}
\end{lstlisting}
\uses{any linear recurrence at a huge index. Fibonacci:
$\begin{pmatrix}1&1\\1&0\end{pmatrix}^n$ gives $F_n$ in
$O(\log n)$. Also counts walks of length $n$ in a graph.}

\subsection{Notable Sums}
\[\sum_{i=1}^{n} i = \frac{n(n+1)}{2} \qquad
  \sum_{i=1}^{n} i^2 = \frac{n(n+1)(2n+1)}{6}\]
\[\sum_{i=1}^{n} i^3 = \left(\frac{n(n+1)}{2}\right)^2 \qquad
  \sum_{i=0}^{n} r^i = \frac{r^{n+1}-1}{r-1}\]
\begin{lstlisting}
ll s1 = n * (n + 1) / 2;
ll s2 = n * (n + 1) * (2 * n + 1) / 6;
\end{lstlisting}
\warn{Compute in \code{ll}: for $n = 10^6$, $n(n+1)$ already overflows
\code{int}.}

\end{multicols}

% =====================================================================
\section{Arrays, Prefix Sums and Searching}
% =====================================================================
\begin{multicols}{2}

\sub{Prefix Sum 1D}{n \text{ build},\; 1 \text{ query}}
\begin{lstlisting}
// pref[i] = sum of the first i elements
vector<ll> pref(n + 1, 0);
for (int i = 0; i < n; i++)
    pref[i + 1] = pref[i] + a[i];

// Inclusive range [l, r], 0-indexed
ll range_sum(int l, int r) { return pref[r + 1] - pref[l]; }
\end{lstlisting}
\tip{Sizing the array $n+1$ and shifting by one removes every
\code{l == 0} special case.}

\subsection{Prefix XOR}
\begin{lstlisting}
vector<int> px(n + 1, 0);
for (int i = 0; i < n; i++) px[i + 1] = px[i] ^ a[i];

int range_xor(int l, int r) { return px[r + 1] ^ px[l]; }
\end{lstlisting}
\uses{XOR is its own inverse, so subtraction becomes another XOR.}

\sub{Prefix Sum 2D}{nm \text{ build},\; 1 \text{ query}}
\begin{lstlisting}
vector<vector<ll>> p(n + 1, vector<ll>(m + 1, 0));
for (int i = 1; i <= n; i++)
    for (int j = 1; j <= m; j++)
        p[i][j] = a[i-1][j-1] + p[i-1][j]
                + p[i][j-1] - p[i-1][j-1];

// Inclusive rectangle, 0-indexed corners
ll rect(int x1, int y1, int x2, int y2) {
    x1++; y1++; x2++; y2++;
    return p[x2][y2] - p[x1-1][y2]
         - p[x2][y1-1] + p[x1-1][y1-1];
}
\end{lstlisting}

\sub{Difference Array 1D}{1 \text{ per update}}
\begin{lstlisting}
vector<ll> d(n + 1, 0);

void add_range(int l, int r, ll v) {   // inclusive
    d[l] += v;
    if (r + 1 <= n) d[r + 1] -= v;
}

// Rebuild the final array once, at the end
ll cur = 0;
for (int i = 0; i < n; i++) { cur += d[i]; a[i] += cur; }
\end{lstlisting}
\uses{many range updates and a single read at the end. It is the exact
mirror of a prefix sum.}

\subsection{Difference Array 2D}
\begin{lstlisting}
vector<vector<ll>> d(n + 2, vector<ll>(m + 2, 0));

void add_rect(int x1, int y1, int x2, int y2, ll v) {
    d[x1][y1]     += v;
    d[x1][y2+1]   -= v;
    d[x2+1][y1]   -= v;
    d[x2+1][y2+1] += v;
}
// Then take a 2D prefix sum of d to get the final grid
\end{lstlisting}

\sub{Binary Search}{\log n}
\begin{lstlisting}
// Classic: find x in a sorted array
int l = 0, r = n - 1, ans = -1;
while (l <= r) {
    int mid = l + (r - l) / 2;     // avoids overflow
    if (a[mid] == x) { ans = mid; break; }
    if (a[mid] < x) l = mid + 1;
    else r = mid - 1;
}
\end{lstlisting}
\warn{\code{(l + r) / 2} overflows when \code{l + r} exceeds
\code{INT\_MAX}. \code{l + (r - l) / 2} never does.}

\sub{Binary Search on the Answer}{\log(\text{range}) \cdot f}
\begin{lstlisting}
// Smallest value that satisfies a monotone predicate
ll lo = 0, hi = 1e18, ans = -1;
while (lo <= hi) {
    ll mid = lo + (hi - lo) / 2;
    if (check(mid)) { ans = mid; hi = mid - 1; }
    else            { lo = mid + 1; }
}
\end{lstlisting}
\uses{"minimise the maximum", "maximise the minimum", and any question
whose answer is monotone: if \code{mid} works then everything above
(or below) works too.}

\subsection{First True Pattern}
\begin{lstlisting}
// Half-open [lo, hi): no ans variable needed
int lo = 0, hi = n;
while (lo < hi) {
    int mid = lo + (hi - lo) / 2;
    if (ok(mid)) hi = mid;
    else         lo = mid + 1;
}
// lo == first index where ok() is true (or n)
\end{lstlisting}
\tip{This shape is much harder to get wrong than the
\code{l <= r} version. Learn one and reuse it everywhere.}

\subsection{Binary Search on Doubles}
\begin{lstlisting}
double lo = 0, hi = 1e9;
for (int iter = 0; iter < 100; iter++) {  // fixed count
    double mid = (lo + hi) / 2;
    if (check(mid)) hi = mid;
    else            lo = mid;
}
\end{lstlisting}
\tip{Iterate a fixed number of times rather than comparing against an
epsilon. 100 halvings is far below any precision you will be asked for,
and it cannot loop forever.}

\sub{lower\_bound / upper\_bound}{\log n}
\begin{lstlisting}
// lower_bound: first element >= x
auto it = lower_bound(all(v), x);
int idx = it - v.begin();
if (it == v.end()) { /* nothing >= x */ }

// upper_bound: first element > x
auto jt = upper_bound(all(v), x);

// Count of x
int cnt = upper_bound(all(v), x) - lower_bound(all(v), x);

// Largest element <= x
auto kt = upper_bound(all(v), x);
if (kt != v.begin()) { --kt; /* *kt <= x */ }
\end{lstlisting}
\warn{Both require the array to be sorted. On unsorted input they return
garbage silently.}

\sub{Ternary Search}{\log(\text{range})}
\begin{lstlisting}
// Maximum of a strictly unimodal function
double lo = 0, hi = 1e9;
for (int iter = 0; iter < 200; iter++) {
    double m1 = lo + (hi - lo) / 3;
    double m2 = hi - (hi - lo) / 3;
    if (f(m1) < f(m2)) lo = m1;
    else               hi = m2;
}
double best = f(lo);
\end{lstlisting}
\uses{unimodal functions --- one peak, no plateau. For a minimum, flip
the comparison.}

\sub{Two Pointers}{n}
\begin{lstlisting}
int l = 0;
ll sum = 0;
for (int r = 0; r < n; r++) {
    sum += a[r];                 // extend the window
    while (sum > K) {            // shrink while invalid
        sum -= a[l];
        l++;
    }
    // [l, r] is now the longest valid window ending at r
    best = max(best, r - l + 1);
}
\end{lstlisting}
\warn{This only works when the predicate is monotone in the window: once
a window is invalid, extending it right cannot make it valid again.}

\sub{Sliding Window Maximum}{n}
\begin{lstlisting}
// Max of every window of size k
deque<int> dq;                   // stores indices
vector<int> res;
for (int i = 0; i < n; i++) {
    while (!dq.empty() && dq.front() <= i - k)
        dq.pop_front();          // drop out-of-window
    while (!dq.empty() && a[dq.back()] <= a[i])
        dq.pop_back();           // drop dominated
    dq.push_back(i);
    if (i >= k - 1) res.push_back(a[dq.front()]);
}
\end{lstlisting}
\tip{The deque stays decreasing, so the front is always the window
maximum. Flip the comparison for a minimum.}

\sub{Kadane: Max Subarray Sum}{n}
\begin{lstlisting}
ll best = LLONG_MIN, cur = 0;
for (int i = 0; i < n; i++) {
    cur = max((ll)a[i], cur + a[i]);
    best = max(best, cur);
}
\end{lstlisting}
\warn{Starting \code{best} at 0 breaks on all-negative arrays. Start it
at \code{LLONG\_MIN} and let the first element set it.}

\sub{LIS}{n \log n}
\begin{lstlisting}
// Length only
int lis_length(const vector<int>& a) {
    vector<int> t;
    for (int x : a) {
        auto it = lower_bound(all(t), x);
        if (it == t.end()) t.push_back(x);
        else *it = x;
    }
    return sz(t);
}
\end{lstlisting}
\warn{\code{t} is \emph{not} the LIS itself, only its length is
meaningful. Reconstruct with parent pointers:}
\begin{lstlisting}
vector<int> lis_reconstruct(const vector<int>& a) {
    int n = sz(a);
    vector<int> t, pos, par(n, -1);
    for (int i = 0; i < n; i++) {
        int j = lower_bound(all(t), a[i]) - t.begin();
        if (j > 0) par[i] = pos[j - 1];
        if (j == sz(t)) { t.push_back(a[i]); pos.push_back(i); }
        else            { t[j] = a[i];       pos[j] = i;       }
    }
    vector<int> res;
    for (int i = pos.back(); i != -1; i = par[i])
        res.push_back(a[i]);
    reverse(all(res));
    return res;
}
\end{lstlisting}
\tip{Swap \code{lower\_bound} for \code{upper\_bound} to get the longest
\emph{non-decreasing} subsequence.}

\sub{Coordinate Compression}{n \log n}
\begin{lstlisting}
vector<int> vals = a;
sort(all(vals));
vals.erase(unique(all(vals)), vals.end());

for (int& x : a)
    x = lower_bound(all(vals), x) - vals.begin();
// vals[x] recovers the original value
\end{lstlisting}
\uses{values up to $10^9$ but only $n \le 10^5$ of them --- compress to
$[0,n)$ so they can index a BIT or segment tree.}

\sub{Counting Inversions}{n \log n}
\begin{lstlisting}
// Pairs i < j with a[i] > a[j], via merge sort
ll merge_count(vector<int>& a, int l, int r) {
    if (r - l <= 1) return 0;
    int m = (l + r) / 2;
    ll cnt = merge_count(a, l, m) + merge_count(a, m, r);

    vector<int> tmp;
    int i = l, j = m;
    while (i < m && j < r) {
        if (a[i] <= a[j]) tmp.push_back(a[i++]);
        else { cnt += m - i; tmp.push_back(a[j++]); }
    }
    while (i < m) tmp.push_back(a[i++]);
    while (j < r) tmp.push_back(a[j++]);
    copy(all(tmp), a.begin() + l);
    return cnt;
}
\end{lstlisting}
\uses{the minimum number of adjacent swaps that sorts the array.}

\sub{Monotonic Stack}{n}
\begin{lstlisting}
// Next strictly greater element to the right
vector<int> nxt(n, -1);
stack<int> st;                  // indices, decreasing values
for (int i = 0; i < n; i++) {
    while (!st.empty() && a[st.top()] < a[i]) {
        nxt[st.top()] = i;
        st.pop();
    }
    st.push(i);
}
\end{lstlisting}
\uses{previous/next greater or smaller, largest rectangle in a
histogram, and the classic "sum of subarray minimums".}

\end{multicols}

% =====================================================================
\section{Bit Manipulation}
% =====================================================================
\begin{multicols}{2}

\subsection{Basic Operations}
\begin{lstlisting}
mask |  (1LL << i)    // set bit i
mask &  (1LL << i)    // test bit i (nonzero if set)
mask ^  (1LL << i)    // toggle bit i
mask & ~(1LL << i)    // clear bit i

(mask >> i) & 1       // test, as a clean 0/1
\end{lstlisting}
\warn{\code{1 << i} is an \code{int}. For $i \ge 31$ this is undefined
behaviour --- write \code{1LL << i} whenever the mask can be wide.}

\subsection{Least Significant Bit}
\begin{lstlisting}
mask & -mask          // isolate the lowest set bit
mask & (mask - 1)     // clear the lowest set bit
mask | (mask + 1)     // set the lowest clear bit
\end{lstlisting}

\subsection{GCC Built-ins}
\begin{lstlisting}
__builtin_popcount(x)    // number of set bits (int)
__builtin_popcountll(x)  // ... (long long)
__builtin_ctz(x)         // trailing zeros -> index of LSB
__builtin_clz(x)         // leading zeros
__builtin_ffs(x)         // 1-based index of LSB, 0 if x == 0
__builtin_parity(x)      // popcount % 2
\end{lstlisting}
\warn{\code{ctz} and \code{clz} are undefined for \code{x == 0}. Guard
them.}
\tip{C++20 offers the portable \code{<bit>} header:
\code{popcount}, \code{countr\_zero}, \code{countl\_zero},
\code{bit\_width}, \code{has\_single\_bit}.}

\subsection{Common Tricks}
\begin{lstlisting}
(1LL << n) - 1            // n low bits all set
n > 0 && !(n & (n - 1))   // is a power of two
31 - __builtin_clz(n)     // floor(log2(n)), int
63 - __builtin_clzll(n)   // floor(log2(n)), long long
1 << (32 - __builtin_clz(n - 1))  // next power of two
x & 1                     // odd?
x >> 1                    // x / 2 for non-negative x
\end{lstlisting}
\warn{Right-shifting a \emph{negative} number is implementation-defined
rounding, not division. Use \code{/ 2} when the sign matters.}

\subsection{XOR Properties}
\begin{itemize}
    \item $a \oplus a = 0$ and $a \oplus 0 = a$.
    \item Commutative and associative --- order never matters.
    \item If $a \oplus b = c$ then $a \oplus c = b$.
    \item XOR of $1..n$ has period 4: $n, 1, n+1, 0$ for
          $n \bmod 4 = 0,1,2,3$.
\end{itemize}
\uses{"every value appears twice except one" --- XOR everything, the
pairs cancel out.}

\sub{Iterate All Subsets}{2^n \cdot n}
\begin{lstlisting}
for (int mask = 0; mask < (1 << n); mask++) {
    for (int j = 0; j < n; j++) {
        if (mask & (1 << j)) {
            // element j belongs to this subset
        }
    }
}
\end{lstlisting}

\sub{Iterate Submasks}{3^n \text{ total}}
\begin{lstlisting}
for (int sub = mask; sub > 0; sub = (sub - 1) & mask) {
    // sub is a nonempty submask of mask
}
// Include the empty submask:
int sub = mask;
while (true) {
    // process sub
    if (sub == 0) break;
    sub = (sub - 1) & mask;
}
\end{lstlisting}
\tip{Summed over every \code{mask}, this loop runs $3^n$ times, not
$4^n$ --- which is what makes subset-sum DP feasible for $n \le 20$.}

\sub{Subsets of Size K (Gosper's Hack)}{\binom{n}{k}}
\begin{lstlisting}
int cur = (1 << k) - 1;
while (cur < (1 << n)) {
    // process cur: exactly k bits set
    int c = cur & -cur;
    int r = cur + c;
    cur = (((r ^ cur) >> 2) / c) | r;
}
\end{lstlisting}

\sub{SOS DP (Sum over Subsets)}{2^n \cdot n}
\begin{lstlisting}
// f[mask] becomes the sum of a[sub] over all sub of mask
vector<ll> f(1 << n);
for (int mask = 0; mask < (1 << n); mask++) f[mask] = a[mask];

for (int i = 0; i < n; i++)
    for (int mask = 0; mask < (1 << n); mask++)
        if (mask & (1 << i))
            f[mask] += f[mask ^ (1 << i)];
\end{lstlisting}
\uses{aggregating over all $2^n$ submasks in $O(2^n n)$ instead of
$O(3^n)$. Reverse the \code{if} to sum over supersets.}

\sub{XOR Linear Basis}{n \log C}
\begin{lstlisting}
// Maximal XOR of any subset of the inserted values
struct Basis {
    static const int B = 60;
    ll b[B + 1] = {};

    bool insert(ll x) {
        for (int i = B; i >= 0; i--) {
            if (!((x >> i) & 1)) continue;
            if (!b[i]) { b[i] = x; return true; }
            x ^= b[i];
        }
        return false;            // x was representable
    }

    ll max_xor() {
        ll res = 0;
        for (int i = B; i >= 0; i--)
            res = max(res, res ^ b[i]);
        return res;
    }
};
\end{lstlisting}
\uses{"maximum XOR of a subset", counting distinct achievable XOR
values ($2^{\text{rank}}$), and checking whether a value is reachable.}

\end{multicols}

% =====================================================================
\section{Strings}
% =====================================================================
\begin{multicols}{2}

\subsection{Reading Input}
\begin{lstlisting}
string s;
getline(cin, s);           // whole line, spaces included
\end{lstlisting}
\warn{After \code{cin >> n}, the newline is still in the buffer and the
next \code{getline} reads an empty string. Clear it first:}
\begin{lstlisting}
int n;
cin >> n;
cin.ignore();              // discard the pending newline
string s;
getline(cin, s);
\end{lstlisting}

\subsection{Tokenising with stringstream}
\begin{lstlisting}
string line = "10 20 30";
stringstream ss(line);
int a, b, c;
ss >> a >> b >> c;

// Unknown count of words
vector<string> words;
string w;
while (ss >> w) words.push_back(w);

// Split on a custom delimiter
while (getline(ss, w, ',')) words.push_back(w);
\end{lstlisting}

\subsection{String Utilities}
\begin{lstlisting}
s.substr(pos, len);        // len is a COUNT, not an index
s.find("abc");             // string::npos if absent
s += 'x';
reverse(all(s));
sort(all(s));              // anagram normal form

int digit = s[i] - '0';
char up = toupper(s[i]);
bool isLetter = isalpha(s[i]);

string t = to_string(42);
int v = stoi(t);
ll  w = stoll(t);
\end{lstlisting}
\warn{\code{s.substr(l, r)} does \emph{not} take a right index. For the
inclusive range $[l,r]$ write \code{s.substr(l, r - l + 1)}.}

\subsection{Character Frequency}
\begin{lstlisting}
array<int, 26> cnt{};
for (char c : s) cnt[c - 'a']++;

// Anagram test
array<int,26> ca{}, cb{};
for (char c : a) ca[c - 'a']++;
for (char c : b) cb[c - 'a']++;
bool anagram = (ca == cb);
\end{lstlisting}

\sub{Z-Function}{n}
\code{z[i]} = length of the longest common prefix of \code{s} and
\code{s[i..]}.
\begin{lstlisting}
vector<int> z_function(const string& s) {
    int n = sz(s);
    vector<int> z(n);
    int l = 0, r = 0;
    for (int i = 1; i < n; i++) {
        if (i <= r) z[i] = min(r - i + 1, z[i - l]);
        while (i + z[i] < n && s[z[i]] == s[i + z[i]]) z[i]++;
        if (i + z[i] - 1 > r) { l = i; r = i + z[i] - 1; }
    }
    return z;
}

vector<int> z_search(const string& text, const string& pat) {
    string s = pat + '#' + text;      // '#' must not occur
    vector<int> z = z_function(s), res;
    int m = sz(pat);
    for (int i = m + 1; i < sz(s); i++)
        if (z[i] == m) res.push_back(i - m - 1);
    return res;
}
\end{lstlisting}
\warn{The separator must be a character that appears in neither string,
or matches can straddle the join.}

\sub{KMP Prefix Function}{n}
\code{pi[i]} = length of the longest proper prefix of \code{s[0..i]}
that is also a suffix of it.
\begin{lstlisting}
vector<int> prefix_function(const string& s) {
    int n = sz(s);
    vector<int> pi(n);
    for (int i = 1; i < n; i++) {
        int j = pi[i - 1];
        while (j > 0 && s[i] != s[j]) j = pi[j - 1];
        if (s[i] == s[j]) j++;
        pi[i] = j;
    }
    return pi;
}

vector<int> kmp_search(const string& text, const string& pat) {
    string s = pat + '#' + text;
    vector<int> pi = prefix_function(s), res;
    int m = sz(pat);
    for (int i = m + 1; i < sz(s); i++)
        if (pi[i] == m) res.push_back(i - 2 * m);
    return res;
}
\end{lstlisting}
\uses{\code{n - pi[n-1]} is the smallest period of the string. It repeats
exactly when that period divides \code{n}.}

\sub{Manacher}{n}
\begin{lstlisting}
// d1[i] = radius of the odd palindrome centred at i
vector<int> manacher_odd(const string& s) {
    int n = sz(s);
    vector<int> d1(n);
    int l = 0, r = -1;
    for (int i = 0; i < n; i++) {
        int k = (i > r) ? 1 : min(d1[l + r - i], r - i + 1);
        while (0 <= i - k && i + k < n && s[i - k] == s[i + k]) k++;
        d1[i] = k;
        if (i + k - 1 > r) { l = i - k + 1; r = i + k - 1; }
    }
    return d1;
}

// d2[i] = radius of the even palindrome between i-1 and i
vector<int> manacher_even(const string& s) {
    int n = sz(s);
    vector<int> d2(n);
    int l = 0, r = -1;
    for (int i = 0; i < n; i++) {
        int k = (i > r) ? 0 : min(d2[l + r - i + 1], r - i + 1);
        while (0 <= i - k - 1 && i + k < n
               && s[i - k - 1] == s[i + k]) k++;
        d2[i] = k;
        if (i + k - 1 > r) { l = i - k; r = i + k - 1; }
    }
    return d2;
}

int longest_palindrome(const string& s) {
    auto d1 = manacher_odd(s), d2 = manacher_even(s);
    int ans = 0;
    for (int x : d1) ans = max(ans, 2 * x - 1);
    for (int x : d2) ans = max(ans, 2 * x);
    return ans;
}
\end{lstlisting}
\uses{the total number of palindromic substrings is
$\sum d1[i] + \sum d2[i]$.}

\sub{String Hashing}{n \text{ build},\; 1 \text{ per substring}}
\begin{lstlisting}
struct Hash {
    static const ll M = (1LL << 61) - 1;   // Mersenne prime
    ll B;                                  // random base
    vector<ll> h, p;

    static ll mul(ll a, ll b) {
        __int128 c = (__int128)a * b;
        ll lo = (ll)(c & M), hi = (ll)(c >> 61);
        lo += hi;
        return lo >= M ? lo - M : lo;
    }

    Hash(const string& s) {
        B = uniform_int_distribution<ll>(256, M - 2)(rng);
        int n = sz(s);
        h.assign(n + 1, 0);
        p.assign(n + 1, 1);
        for (int i = 0; i < n; i++) {
            h[i + 1] = (mul(h[i], B) + s[i]) % M;
            p[i + 1] = mul(p[i], B);
        }
    }

    // Hash of the inclusive range [l, r]
    ll get(int l, int r) {
        ll x = (h[r + 1] - mul(h[l], p[r - l + 1])) % M;
        return x < 0 ? x + M : x;
    }
};
\end{lstlisting}
\uses{$O(1)$ substring comparison, so binary searching the longest
common prefix of two suffixes costs $O(\log n)$.}
\warn{Randomise the base. A fixed base with a small modulus is trivially
hacked on Codeforces. Modulo $2^{61}-1$ with a random base is safe in
practice.}

\sub{Trie}{|s| \text{ per operation}}
\begin{lstlisting}
struct Trie {
    static const int A = 26;
    vector<array<int, A>> nxt;
    vector<int> cnt;                  // words ending here

    Trie() { new_node(); }

    int new_node() {
        nxt.push_back({});
        nxt.back().fill(-1);
        cnt.push_back(0);
        return sz(nxt) - 1;
    }

    void insert(const string& s) {
        int v = 0;
        for (char ch : s) {
            int c = ch - 'a';
            if (nxt[v][c] == -1) nxt[v][c] = new_node();
            v = nxt[v][c];
        }
        cnt[v]++;
    }

    bool contains(const string& s) {
        int v = 0;
        for (char ch : s) {
            int c = ch - 'a';
            if (nxt[v][c] == -1) return false;
            v = nxt[v][c];
        }
        return cnt[v] > 0;
    }
};
\end{lstlisting}
\uses{prefix queries, autocomplete, and --- with \code{A = 2} over the
bits of each number --- maximum XOR pair queries.}

\end{multicols}

% =====================================================================
\section{Graphs: Traversal and Ordering}
% =====================================================================
\begin{multicols}{2}

\subsection{Representations}
\begin{lstlisting}
const int MAXN = 2e5 + 5;

// Adjacency list: the default choice
vector<int> adj[MAXN];
adj[u].push_back(v);
adj[v].push_back(u);              // omit if directed

// Weighted
vector<pii> wadj[MAXN];           // {neighbour, weight}
wadj[u].push_back({v, w});

// Edge list: for Kruskal and Bellman-Ford
vector<tuple<int,int,int>> edges; // {w, u, v}
\end{lstlisting}

\subsection{Graph Identities}
\begin{itemize}
    \item Tree on $n$ nodes: exactly $n - 1$ edges, connected, acyclic.
    \item Complete graph $K_n$: $\frac{n(n-1)}{2}$ edges.
    \item Handshake lemma: $\sum_{v} \deg(v) = 2|E|$.
    \item A connected graph has an Euler circuit iff every degree is
          even; an Euler path iff exactly two degrees are odd.
    \item A graph is bipartite iff it has no odd-length cycle.
\end{itemize}

\sub{DFS}{V + E}
\begin{lstlisting}
bool vis[MAXN];

void dfs(int u) {
    vis[u] = true;
    for (int v : adj[u])
        if (!vis[v]) dfs(v);
}
\end{lstlisting}
\warn{Recursion depth equals the longest path. For $n = 10^5$ on a path
graph this overflows the default stack. Use the iterative form below.}

\subsection{Iterative DFS}
\begin{lstlisting}
void dfs_iter(int s) {
    vector<int> st{s};
    vis[s] = true;
    while (!st.empty()) {
        int u = st.back(); st.pop_back();
        for (int v : adj[u])
            if (!vis[v]) { vis[v] = true; st.push_back(v); }
    }
}
\end{lstlisting}
\tip{Mark nodes visited when you \emph{push} them, not when you pop, or
the same node lands on the stack many times.}

\sub{BFS and Distances}{V + E}
\begin{lstlisting}
vector<int> bfs(int s, int n) {
    vector<int> dist(n + 1, -1);
    queue<int> q;
    q.push(s);
    dist[s] = 0;
    while (!q.empty()) {
        int u = q.front(); q.pop();
        for (int v : adj[u])
            if (dist[v] == -1) {
                dist[v] = dist[u] + 1;
                q.push(v);
            }
    }
    return dist;
}
\end{lstlisting}
\uses{shortest paths on \emph{unweighted} graphs. \code{dist[v] == -1}
doubles as the visited flag.}

\subsection{Multi-Source BFS}
\begin{lstlisting}
// Distance from the NEAREST source, in one pass
queue<int> q;
for (int s : sources) { dist[s] = 0; q.push(s); }
// ... identical BFS loop from here
\end{lstlisting}
\uses{"distance to the closest fire / gate / water cell". Seeding every
source at once beats running one BFS per source.}

\sub{0-1 BFS}{V + E}
\begin{lstlisting}
// Shortest path when every edge weight is 0 or 1
deque<int> dq;
vector<int> dist(n, INT_MAX);
dist[s] = 0;
dq.push_front(s);
while (!dq.empty()) {
    int u = dq.front(); dq.pop_front();
    for (auto [v, w] : wadj[u]) {
        if (dist[u] + w < dist[v]) {
            dist[v] = dist[u] + w;
            if (w == 0) dq.push_front(v);
            else        dq.push_back(v);
        }
    }
}
\end{lstlisting}
\uses{a Dijkstra replacement without the $\log$ factor when the only
weights are 0 and 1.}

\subsection{Grid Traversal}
\begin{lstlisting}
int dx[4] = {-1, 1, 0, 0};
int dy[4] = { 0, 0,-1, 1};

// 8 directions, including diagonals
int ex[8] = {-1,-1,-1, 0, 0, 1, 1, 1};
int ey[8] = {-1, 0, 1,-1, 1,-1, 0, 1};

for (int k = 0; k < 4; k++) {
    int nx = x + dx[k], ny = y + dy[k];
    if (nx < 0 || nx >= R || ny < 0 || ny >= C) continue;
    if (grid[nx][ny] == '#') continue;
    // ...
}
\end{lstlisting}
\tip{Flatten a cell to \code{x * C + y} when you need it as a single
integer id (for a DSU, or a \code{dist} array).}

\sub{Connected Components}{V + E}
\begin{lstlisting}
int components = 0;
for (int i = 1; i <= n; i++)
    if (!vis[i]) { components++; dfs(i); }
\end{lstlisting}

\sub{Bipartite Check}{V + E}
\begin{lstlisting}
vector<int> color(MAXN, -1);

bool bipartite(int n) {
    for (int s = 1; s <= n; s++) {
        if (color[s] != -1) continue;
        color[s] = 0;
        queue<int> q;
        q.push(s);
        while (!q.empty()) {
            int u = q.front(); q.pop();
            for (int v : adj[u]) {
                if (color[v] == -1) {
                    color[v] = color[u] ^ 1;
                    q.push(v);
                } else if (color[v] == color[u]) {
                    return false;
                }
            }
        }
    }
    return true;
}
\end{lstlisting}
\warn{Loop over every start node. Checking only node 1 misses
disconnected components.}

\sub{Cycle in an Undirected Graph}{V + E}
\begin{lstlisting}
bool dfs_cycle(int u, int parent) {
    vis[u] = true;
    for (int v : adj[u]) {
        if (!vis[v]) {
            if (dfs_cycle(v, u)) return true;
        } else if (v != parent) {
            return true;
        }
    }
    return false;
}
\end{lstlisting}
\warn{With multi-edges, comparing to \code{parent} by node id is wrong
--- track the incoming edge id instead.}

\sub{Cycle in a Directed Graph}{V + E}
\begin{lstlisting}
// 0 = unvisited, 1 = on the current path, 2 = done
int state[MAXN];

bool dfs_cycle_dir(int u) {
    state[u] = 1;
    for (int v : adj[u]) {
        if (state[v] == 1) return true;      // back edge
        if (state[v] == 0 && dfs_cycle_dir(v)) return true;
    }
    state[u] = 2;
    return false;
}
\end{lstlisting}
\warn{State 2 matters: revisiting a \emph{finished} node is a cross
edge, not a cycle. A plain \code{vis} array reports false cycles.}

\sub{Topological Sort (Kahn)}{V + E}
\begin{lstlisting}
vector<int> topo_sort(int n) {
    vector<int> indeg(n + 1, 0), order;
    for (int u = 1; u <= n; u++)
        for (int v : adj[u]) indeg[v]++;

    queue<int> q;                 // priority_queue for
    for (int i = 1; i <= n; i++)  // lexicographically smallest
        if (indeg[i] == 0) q.push(i);

    while (!q.empty()) {
        int u = q.front(); q.pop();
        order.push_back(u);
        for (int v : adj[u])
            if (--indeg[v] == 0) q.push(v);
    }
    if (sz(order) != n) return {};   // a cycle exists
    return order;
}
\end{lstlisting}
\uses{a short result means the graph has a cycle --- this doubles as
cycle detection. Longest path on a DAG is a DP over this order.}

\subsection{Topological Sort (DFS)}
\begin{lstlisting}
void topo_dfs(int u, vector<int>& order) {
    vis[u] = true;
    for (int v : adj[u]) if (!vis[v]) topo_dfs(v, order);
    order.push_back(u);            // post-order
}
// Reverse `order` at the end
\end{lstlisting}

\end{multicols}

% =====================================================================
\section{Shortest Paths, MST and DSU}
% =====================================================================
\begin{multicols}{2}

\sub{Dijkstra}{E \log V}
\begin{lstlisting}
vector<ll> dijkstra(int s, int n) {
    vector<ll> dist(n + 1, INF);
    priority_queue<pll, vector<pll>, greater<pll>> pq;
    dist[s] = 0;
    pq.push({0, s});

    while (!pq.empty()) {
        auto [d, u] = pq.top(); pq.pop();
        if (d > dist[u]) continue;        // stale entry

        for (auto [v, w] : wadj[u]) {
            if (dist[u] + w < dist[v]) {
                dist[v] = dist[u] + w;
                pq.push({dist[v], v});
            }
        }
    }
    return dist;
}
\end{lstlisting}
\warn{Fails with negative weights. The \code{if (d > dist[u]) continue}
line is not optional --- without it, stale queue entries reprocess
nodes and the complexity degrades.}

\subsection{Path Reconstruction}
\begin{lstlisting}
vector<int> par(n + 1, -1);
// inside the relaxation:
//   par[v] = u;

vector<int> path;
for (int cur = t; cur != -1; cur = par[cur])
    path.push_back(cur);
reverse(all(path));
\end{lstlisting}

\sub{Bellman-Ford}{VE}
\begin{lstlisting}
// Handles negative weights and detects negative cycles
vector<ll> bellman_ford(int s, int n, bool& has_neg_cycle) {
    vector<ll> dist(n + 1, INF);
    dist[s] = 0;
    for (int it = 0; it < n - 1; it++)
        for (auto [w, u, v] : edges)
            if (dist[u] < INF && dist[u] + w < dist[v])
                dist[v] = dist[u] + w;

    has_neg_cycle = false;
    for (auto [w, u, v] : edges)          // one extra pass
        if (dist[u] < INF && dist[u] + w < dist[v])
            has_neg_cycle = true;
    return dist;
}
\end{lstlisting}
\uses{a distance that still improves on pass $n$ can only mean a
negative cycle.}

\sub{Floyd-Warshall}{V^3}
\begin{lstlisting}
// d[i][j] = weight of edge i->j, INF if none, 0 on diagonal
for (int k = 1; k <= n; k++)
    for (int i = 1; i <= n; i++) {
        if (d[i][k] == INF) continue;     // big speedup
        for (int j = 1; j <= n; j++)
            if (d[i][k] + d[k][j] < d[i][j])
                d[i][j] = d[i][k] + d[k][j];
    }

// d[i][i] < 0  =>  i lies on a negative cycle
\end{lstlisting}
\warn{\code{k} must be the outermost loop. Any other order silently
computes the wrong answer. Practical up to $n \approx 500$.}
\tip{With \code{d} as booleans and \code{|=}/\code{\&} instead of
\code{+}/\code{min}, the same loop is transitive closure
(reachability).}

\sub{DSU / Union-Find}{\alpha(n)}
\begin{lstlisting}
struct DSU {
    vector<int> parent, sz;
    int components;

    DSU(int n) : parent(n + 1), sz(n + 1, 1), components(n) {
        iota(all(parent), 0);
    }

    int find(int x) {                     // path compression
        return parent[x] == x
             ? x
             : parent[x] = find(parent[x]);
    }

    bool unite(int a, int b) {            // union by size
        a = find(a); b = find(b);
        if (a == b) return false;
        if (sz[a] < sz[b]) swap(a, b);
        parent[b] = a;
        sz[a] += sz[b];
        components--;
        return true;
    }

    bool same(int a, int b) { return find(a) == find(b); }
    int  size(int x)        { return sz[find(x)];        }
};
\end{lstlisting}
\uses{connectivity under edge insertions, Kruskal, and counting
components online. $\alpha(n) \le 4$ for any realistic $n$ --- treat it
as constant.}
\warn{DSU cannot delete edges. If the problem removes edges, process
the queries in reverse.}

\sub{Kruskal MST}{E \log E}
\begin{lstlisting}
ll kruskal(int n, vector<tuple<int,int,int>>& edges) {
    sort(all(edges));                     // by weight
    DSU dsu(n);
    ll total = 0;
    int used = 0;
    for (auto [w, u, v] : edges) {
        if (dsu.unite(u, v)) {
            total += w;
            if (++used == n - 1) break;
        }
    }
    return used == n - 1 ? total : -1;    // -1: disconnected
}
\end{lstlisting}
\tip{Sort descending for a maximum spanning tree. For the "minimise the
largest edge on a path" family, the MST is the answer.}

\sub{Prim MST}{E \log V}
\begin{lstlisting}
// Better than Kruskal on dense graphs
ll prim(int n) {
    vector<bool> in(n + 1, false);
    priority_queue<pll, vector<pll>, greater<pll>> pq;
    pq.push({0, 1});
    ll total = 0;
    int used = 0;
    while (!pq.empty() && used < n) {
        auto [w, u] = pq.top(); pq.pop();
        if (in[u]) continue;
        in[u] = true;
        total += w;
        used++;
        for (auto [v, c] : wadj[u])
            if (!in[v]) pq.push({c, v});
    }
    return used == n ? total : -1;
}
\end{lstlisting}

\end{multicols}

% =====================================================================
\section{Trees and Advanced Graphs}
% =====================================================================
\begin{multicols}{2}

\sub{Rooting: Parent, Depth, Subtree Size}{n}
\begin{lstlisting}
vector<int> par, dep, sub;

void root_tree(int u, int p) {
    par[u] = p;
    sub[u] = 1;
    for (int v : adj[u]) {
        if (v == p) continue;
        dep[v] = dep[u] + 1;
        root_tree(v, u);
        sub[u] += sub[v];
    }
}
// call: dep[1] = 0; root_tree(1, 0);
\end{lstlisting}
\tip{Skipping the parent is all a tree DFS needs --- no visited array.}

\sub{Tree Diameter}{n}
\begin{lstlisting}
// Two BFS passes. Works on unweighted trees.
// 1. BFS from any node -> the farthest node a
// 2. BFS from a        -> the farthest node b
// dist(a, b) is the diameter
int farthest(int s, int n) {
    auto d = bfs(s, n);
    int best = s;
    for (int i = 1; i <= n; i++)
        if (d[i] > d[best]) best = i;
    return best;
}
\end{lstlisting}
\warn{The two-BFS trick is only valid on trees (and on graphs with
non-negative weights using Dijkstra). It is wrong on general graphs.}

\sub{LCA with Binary Lifting}{n \log n \text{ build},\; \log n \text{ query}}
\begin{lstlisting}
const int LOG = 18;              // 2^18 > 2*10^5
int up[MAXN][LOG];

void build_lca(int n) {
    for (int i = 1; i <= n; i++) up[i][0] = par[i];
    for (int j = 1; j < LOG; j++)
        for (int i = 1; i <= n; i++)
            up[i][j] = up[ up[i][j-1] ][j-1];
}

int lca(int u, int v) {
    if (dep[u] < dep[v]) swap(u, v);
    int diff = dep[u] - dep[v];
    for (int j = 0; j < LOG; j++)          // lift u up
        if (diff & (1 << j)) u = up[u][j];
    if (u == v) return u;

    for (int j = LOG - 1; j >= 0; j--)     // lift together
        if (up[u][j] != up[v][j]) {
            u = up[u][j];
            v = up[v][j];
        }
    return up[u][0];
}

int dist(int u, int v) {
    return dep[u] + dep[v] - 2 * dep[lca(u, v)];
}

int kth_ancestor(int u, int k) {
    for (int j = 0; j < LOG && u; j++)
        if (k & (1 << j)) u = up[u][j];
    return u;
}
\end{lstlisting}
\warn{Set \code{par[root] = root} (or 0 with \code{up[0][*] = 0}) so
lifting past the root is harmless.}

\sub{Euler Tour: Subtree as a Range}{n}
\begin{lstlisting}
int timer = 0;
vector<int> tin(MAXN), tout(MAXN);

void euler(int u, int p) {
    tin[u] = timer++;
    for (int v : adj[u]) if (v != p) euler(v, u);
    tout[u] = timer;               // half-open [tin, tout)
}

// u is an ancestor of v
bool is_anc(int u, int v) {
    return tin[u] <= tin[v] && tout[v] <= tout[u];
}
\end{lstlisting}
\uses{the subtree of \code{u} is exactly the contiguous range
$[tin_u, tout_u)$. Subtree sum and subtree update become an ordinary BIT
or segment tree range operation.}

\sub{DP on Trees}{n}
\begin{lstlisting}
// Example: largest independent set in a tree
// dp[u][0] = u not taken, dp[u][1] = u taken
ll dp[MAXN][2];

void tree_dp(int u, int p) {
    dp[u][0] = 0;
    dp[u][1] = 1;
    for (int v : adj[u]) {
        if (v == p) continue;
        tree_dp(v, u);
        dp[u][0] += max(dp[v][0], dp[v][1]);
        dp[u][1] += dp[v][0];
    }
}
// answer: max(dp[root][0], dp[root][1])
\end{lstlisting}

\sub{Bridges}{V + E}
An edge whose removal disconnects the graph.
\begin{lstlisting}
vector<int> tin_b(MAXN, -1), low(MAXN);
vector<pii> bridges;
int timer_b = 0;

void dfs_bridge(int u, int pe) {   // pe = incoming edge id
    tin_b[u] = low[u] = timer_b++;
    for (auto [v, id] : adj_id[u]) {
        if (id == pe) continue;
        if (tin_b[v] != -1) {
            low[u] = min(low[u], tin_b[v]);
        } else {
            dfs_bridge(v, id);
            low[u] = min(low[u], low[v]);
            if (low[v] > tin_b[u]) bridges.push_back({u, v});
        }
    }
}
\end{lstlisting}
\warn{Skip by \emph{edge id}, not by parent node --- otherwise parallel
edges are misreported as bridges. \code{adj\_id[u]} stores
\code{\{neighbour, edge\_id\}}.}

\sub{Articulation Points}{V + E}
\begin{lstlisting}
vector<bool> is_art(MAXN, false);

void dfs_art(int u, int p) {
    tin_b[u] = low[u] = timer_b++;
    int children = 0;
    for (int v : adj[u]) {
        if (v == p) continue;
        if (tin_b[v] != -1) {
            low[u] = min(low[u], tin_b[v]);
        } else {
            dfs_art(v, u);
            low[u] = min(low[u], low[v]);
            if (low[v] >= tin_b[u] && p != -1) is_art[u] = true;
            children++;
        }
    }
    if (p == -1 && children > 1) is_art[u] = true;
}
\end{lstlisting}
\warn{The root is special: it is an articulation point only when it has
more than one DFS child.}

\sub{SCC (Kosaraju)}{V + E}
\begin{lstlisting}
vector<int> order_k, comp;
vector<int> radj[MAXN];             // reversed graph

void dfs1(int u) {
    vis[u] = true;
    for (int v : adj[u]) if (!vis[v]) dfs1(v);
    order_k.push_back(u);
}

void dfs2(int u, int c) {
    comp[u] = c;
    for (int v : radj[u]) if (comp[v] == -1) dfs2(v, c);
}

int kosaraju(int n) {
    comp.assign(n + 1, -1);
    for (int i = 1; i <= n; i++) if (!vis[i]) dfs1(i);
    reverse(all(order_k));
    int c = 0;
    for (int u : order_k)
        if (comp[u] == -1) dfs2(u, c++);
    return c;                       // number of SCCs
}
\end{lstlisting}
\uses{contracting each SCC to a single node yields a DAG (the condensation),
which you can then topologically sort or DP over.}

\end{multicols}

% =====================================================================
\section{Dynamic Programming}
% =====================================================================
\begin{multicols}{2}

\subsection{How to Build a DP}
\begin{enumerate}
    \item \textbf{State} --- what exactly does \code{dp[i][j]} mean?
          Write it down as a sentence first.
    \item \textbf{Transition} --- how do smaller states combine into it?
    \item \textbf{Base case} --- the smallest state, filled directly.
    \item \textbf{Order} --- iterate so every dependency is already
          computed.
    \item \textbf{Answer} --- which state (or which max/sum over states)
          is the result?
\end{enumerate}
\tip{Total cost $=$ number of states $\times$ cost per transition.
Check that against the constraints before writing any code.}

\sub{0/1 Knapsack}{nW}
\begin{lstlisting}
// Each item used at most once
vector<ll> dp(W + 1, 0);
for (int i = 0; i < n; i++)
    for (int w = W; w >= wt[i]; w--)      // DESCENDING
        dp[w] = max(dp[w], dp[w - wt[i]] + val[i]);
// dp[W] = best value within capacity W
\end{lstlisting}
\warn{The descending loop is what enforces "at most once". Ascending
turns it into the unbounded version.}

\sub{Unbounded Knapsack / Coin Change}{nW}
\begin{lstlisting}
// Minimum number of coins summing to W
vector<int> dp(W + 1, IINF);
dp[0] = 0;
for (int c : coins)
    for (int w = c; w <= W; w++)          // ASCENDING
        if (dp[w - c] != IINF)
            dp[w] = min(dp[w], dp[w - c] + 1);

// Number of distinct combinations (order ignored)
vector<ll> ways(W + 1, 0);
ways[0] = 1;
for (int c : coins)
    for (int w = c; w <= W; w++)
        ways[w] += ways[w - c];
\end{lstlisting}
\warn{Loop order decides the meaning. Coins outside, sums inside counts
\emph{combinations}; sums outside, coins inside counts
\emph{permutations}.}

\sub{LCS}{nm}
\begin{lstlisting}
vector<vector<int>> dp(n + 1, vector<int>(m + 1, 0));
for (int i = 1; i <= n; i++)
    for (int j = 1; j <= m; j++)
        dp[i][j] = (a[i-1] == b[j-1])
                 ? dp[i-1][j-1] + 1
                 : max(dp[i-1][j], dp[i][j-1]);
\end{lstlisting}

\sub{Edit Distance}{nm}
\begin{lstlisting}
vector<vector<int>> dp(n + 1, vector<int>(m + 1, 0));
for (int i = 0; i <= n; i++) dp[i][0] = i;   // deletions
for (int j = 0; j <= m; j++) dp[0][j] = j;   // insertions

for (int i = 1; i <= n; i++)
    for (int j = 1; j <= m; j++)
        dp[i][j] = (a[i-1] == b[j-1])
            ? dp[i-1][j-1]
            : 1 + min({dp[i-1][j],      // delete
                       dp[i][j-1],      // insert
                       dp[i-1][j-1]});  // replace
\end{lstlisting}

\sub{Bitmask DP (TSP)}{2^n n^2}
\begin{lstlisting}
// dp[mask][i] = min cost visiting `mask`, standing at i
vector<vector<ll>> dp(1 << n, vector<ll>(n, INF));
dp[1][0] = 0;                         // start at node 0

for (int mask = 1; mask < (1 << n); mask++)
    for (int i = 0; i < n; i++) {
        if (dp[mask][i] == INF) continue;
        if (!(mask & (1 << i))) continue;
        for (int j = 0; j < n; j++) {
            if (mask & (1 << j)) continue;
            int nm = mask | (1 << j);
            dp[nm][j] = min(dp[nm][j], dp[mask][i] + cost[i][j]);
        }
    }

ll ans = INF;
for (int i = 1; i < n; i++)
    ans = min(ans, dp[(1 << n) - 1][i] + cost[i][0]);
\end{lstlisting}
\uses{$n \le 20$ is the giveaway that a bitmask DP is intended.}

\sub{Digit DP}{|d| \cdot 10 \cdot \text{states}}
\begin{lstlisting}
string num;
ll memo[20][2][200];
bool seen[20][2][200];

// Counts values in [0, num] with some digit property
ll go(int pos, bool tight, int acc) {
    if (pos == sz(num)) return good(acc) ? 1 : 0;
    if (!tight && seen[pos][tight][acc])
        return memo[pos][tight][acc];

    int lim = tight ? num[pos] - '0' : 9;
    ll res = 0;
    for (int d = 0; d <= lim; d++)
        res += go(pos + 1, tight && (d == lim), acc + d);

    if (!tight) {
        seen[pos][tight][acc] = true;
        memo[pos][tight][acc] = res;
    }
    return res;
}
// Range [L, R]: go(R) - go(L - 1)
\end{lstlisting}
\warn{Only memoise the \code{!tight} states --- tight states depend on
the specific prefix and are not reusable.}

\sub{Longest Path on a DAG}{V + E}
\begin{lstlisting}
// Process nodes in topological order
vector<ll> dp(n + 1, 0);
for (int u : topo_order)
    for (int v : adj[u])
        dp[v] = max(dp[v], dp[u] + 1);
\end{lstlisting}
\uses{longest path is NP-hard on general graphs but linear on a DAG,
because the topological order removes all cyclic dependencies.}

\subsection{Memory Optimisation}
\begin{lstlisting}
// When dp[i] only reads dp[i-1], keep two rows
vector<int> prev(m + 1, 0), cur(m + 1, 0);
for (int i = 1; i <= n; i++) {
    for (int j = 1; j <= m; j++) cur[j] = /* ... */;
    swap(prev, cur);
}
\end{lstlisting}
\tip{$10^5 \times 10^5$ \code{int}s is 40 GB; two rows is 800 KB.
Rolling arrays are often the difference between MLE and AC.}

\end{multicols}

% =====================================================================
\section{Advanced Data Structures}
% =====================================================================
\begin{multicols}{2}

\sub{Fenwick Tree (BIT)}{\log n}
\begin{lstlisting}
struct Fenwick {
    int n;
    vector<ll> bit;

    Fenwick(int n) : n(n), bit(n + 1, 0) {}

    void add(int i, ll v) {          // 1-indexed
        for (; i <= n; i += i & -i) bit[i] += v;
    }

    ll sum(int i) {                  // prefix [1, i]
        ll r = 0;
        for (; i > 0; i -= i & -i) r += bit[i];
        return r;
    }

    ll range(int l, int r) { return sum(r) - sum(l - 1); }

    // Smallest i with prefix sum >= target, O(log n)
    int lower_bound(ll target) {
        int pos = 0;
        for (int pw = 1 << __lg(n); pw; pw >>= 1)
            if (pos + pw <= n && bit[pos + pw] < target) {
                pos += pw;
                target -= bit[pos];
            }
        return pos + 1;
    }
};
\end{lstlisting}
\warn{A BIT is 1-indexed. Feeding it index 0 makes \code{add} loop
forever, because \code{0 \& -0 == 0}.}
\tip{Smaller and roughly twice as fast as a segment tree. Reach for a
BIT whenever prefix sums are enough.}

\sub{Fenwick: Range Update, Point Query}{\log n}
\begin{lstlisting}
// Store the difference array inside the BIT
void add_range(Fenwick& f, int l, int r, ll v) {
    f.add(l, v);
    f.add(r + 1, -v);
}
ll point_query(Fenwick& f, int i) { return f.sum(i); }
\end{lstlisting}

\sub{2D Fenwick}{\log^2 n}
\begin{lstlisting}
ll bit2[N + 1][M + 1];

void add2(int x, int y, ll v) {
    for (int i = x; i <= N; i += i & -i)
        for (int j = y; j <= M; j += j & -j)
            bit2[i][j] += v;
}

ll sum2(int x, int y) {
    ll r = 0;
    for (int i = x; i > 0; i -= i & -i)
        for (int j = y; j > 0; j -= j & -j)
            r += bit2[i][j];
    return r;
}

ll rect2(int x1, int y1, int x2, int y2) {
    return sum2(x2, y2) - sum2(x1 - 1, y2)
         - sum2(x2, y1 - 1) + sum2(x1 - 1, y1 - 1);
}
\end{lstlisting}

\sub{Segment Tree (Recursive)}{\log n}
\begin{lstlisting}
const int MAXN = 2e5 + 5;
ll t[4 * MAXN];
int a[MAXN];

void build(int v, int tl, int tr) {
    if (tl == tr) { t[v] = a[tl]; return; }
    int tm = (tl + tr) / 2;
    build(2*v, tl, tm);
    build(2*v + 1, tm + 1, tr);
    t[v] = t[2*v] + t[2*v + 1];
}

ll query(int v, int tl, int tr, int l, int r) {
    if (l > r) return 0;                 // neutral element
    if (l == tl && r == tr) return t[v];
    int tm = (tl + tr) / 2;
    return query(2*v, tl, tm, l, min(r, tm))
         + query(2*v + 1, tm + 1, tr, max(l, tm + 1), r);
}

void update(int v, int tl, int tr, int pos, ll val) {
    if (tl == tr) { t[v] = val; return; }
    int tm = (tl + tr) / 2;
    if (pos <= tm) update(2*v, tl, tm, pos, val);
    else           update(2*v + 1, tm + 1, tr, pos, val);
    t[v] = t[2*v] + t[2*v + 1];
}
\end{lstlisting}
\tip{For min/max, swap \code{+} for \code{min}/\code{max} and change the
neutral element to \code{INF}/\code{-INF}. Size the array $4n$, not
$2n$.}

\sub{Segment Tree (Iterative)}{\log n}
\begin{lstlisting}
// Shorter, ~2x faster, point update + range query
int N;
ll t[2 * MAXN];

void build_it(int n) {
    N = n;
    for (int i = N - 1; i > 0; i--) t[i] = t[2*i] + t[2*i+1];
}

void update_it(int p, ll val) {
    for (t[p += N] = val; p > 1; p >>= 1)
        t[p >> 1] = t[p] + t[p ^ 1];
}

ll query_it(int l, int r) {          // inclusive [l, r]
    ll res = 0;
    for (l += N, r += N + 1; l < r; l >>= 1, r >>= 1) {
        if (l & 1) res += t[l++];
        if (r & 1) res += t[--r];
    }
    return res;
}
\end{lstlisting}
\warn{Leaves live at \code{t[N .. 2N-1]}. Fill them before calling
\code{build\_it}.}

\sub{Segment Tree with Lazy Propagation}{\log n}
\begin{lstlisting}
ll t[4 * MAXN], lz[4 * MAXN];

void push(int v, int tl, int tr) {
    if (!lz[v]) return;
    t[v] += lz[v] * (tr - tl + 1);
    if (tl != tr) {
        lz[2*v]     += lz[v];
        lz[2*v + 1] += lz[v];
    }
    lz[v] = 0;
}

void update(int v, int tl, int tr, int l, int r, ll val) {
    push(v, tl, tr);
    if (l > tr || tl > r) return;
    if (l <= tl && tr <= r) {
        lz[v] += val;
        push(v, tl, tr);
        return;
    }
    int tm = (tl + tr) / 2;
    update(2*v, tl, tm, l, r, val);
    update(2*v + 1, tm + 1, tr, l, r, val);
    t[v] = t[2*v] + t[2*v + 1];
}

ll query(int v, int tl, int tr, int l, int r) {
    push(v, tl, tr);
    if (l > tr || tl > r) return 0;
    if (l <= tl && tr <= r) return t[v];
    int tm = (tl + tr) / 2;
    return query(2*v, tl, tm, l, r)
         + query(2*v + 1, tm + 1, tr, l, r);
}
\end{lstlisting}
\warn{\code{push} must run at the top of \emph{every} recursive call,
including queries. A forgotten push is the classic lazy-propagation bug.}

\sub{Sparse Table (RMQ)}{n \log n \text{ build},\; 1 \text{ query}}
\begin{lstlisting}
int sp[LOG][MAXN], lg[MAXN];

void build_sparse(int n) {
    lg[1] = 0;
    for (int i = 2; i <= n; i++) lg[i] = lg[i / 2] + 1;

    for (int i = 0; i < n; i++) sp[0][i] = a[i];
    for (int j = 1; j < LOG; j++)
        for (int i = 0; i + (1 << j) <= n; i++)
            sp[j][i] = min(sp[j-1][i],
                           sp[j-1][i + (1 << (j-1))]);
}

int query_min(int l, int r) {        // inclusive
    int j = lg[r - l + 1];
    return min(sp[j][l], sp[j][r - (1 << j) + 1]);
}
\end{lstlisting}
\warn{Only valid for \emph{idempotent} operations --- min, max, gcd ---
where overlapping ranges are harmless. It does not work for sums, and it
does not support updates.}

\sub{PBDS Ordered Set}{\log n}
\begin{lstlisting}
#include <ext/pb_ds/assoc_container.hpp>
#include <ext/pb_ds/tree_policy.hpp>
using namespace __gnu_pbds;

typedef tree<int, null_type, less<int>, rb_tree_tag,
             tree_order_statistics_node_update> ordered_set;

ordered_set s;
s.insert(1);
s.insert(9);

// k-th smallest, 0-indexed
auto it = s.find_by_order(0);
// how many elements are strictly smaller than x
int cnt = s.order_of_key(5);
\end{lstlisting}
\warn{GCC only. It rejects duplicates like a \code{set}; for a multiset
use \code{less\_equal<int>} as the comparator --- but then
\code{erase(value)} stops working and you must erase by iterator.}

\end{multicols}

% =====================================================================
\section{Contest Techniques and Pitfalls}
% =====================================================================
\begin{multicols}{2}

\sub{Meet in the Middle}{2^{n/2} \cdot n}
\begin{lstlisting}
// Subset sums for n <= 40: split into two halves
vector<ll> gen(vector<ll>& v) {
    int k = sz(v);
    vector<ll> sums;
    for (int mask = 0; mask < (1 << k); mask++) {
        ll s = 0;
        for (int i = 0; i < k; i++)
            if (mask & (1 << i)) s += v[i];
        sums.push_back(s);
    }
    return sums;
}
// Sort one half, binary search it for each element
// of the other half.
\end{lstlisting}
\uses{$n \le 40$ where $2^n$ is hopeless but $2^{20}$ is instant.}

\subsection{Interactive Problems}
\begin{lstlisting}
// Flush after every query or you will deadlock
cout << "? " << x << endl;   // endl flushes
cin >> response;

cout << "! " << answer << endl;
\end{lstlisting}
\warn{This is the one place where \code{endl} is correct and
\code{'\textbackslash n'} is a bug.}

\subsection{Stress Testing}
When a solution fails on a hidden test, generate your own:
\begin{lstlisting}
# gen.py writes a small random test to stdout
for i in $(seq 1 1000); do
    python gen.py $i > in.txt
    ./fast  < in.txt > out1.txt
    ./brute < in.txt > out2.txt
    if ! diff -q out1.txt out2.txt > /dev/null; then
        echo "Mismatch on test $i"
        cat in.txt
        break
    fi
done
\end{lstlisting}
\tip{Generate \emph{small} tests ($n \le 8$). A minimal counterexample
you can read beats a 200-line one you cannot.}

\subsection{Greedy: Proving It Works}
\begin{itemize}
    \item \textbf{Exchange argument} --- show any optimal solution can be
          rewritten into the greedy one without getting worse.
    \item \textbf{Stays ahead} --- show the greedy prefix is at least as
          good as any other prefix at every step.
    \item If neither works, it is probably a DP.
\end{itemize}
\tip{Cheapest check: brute-force small cases and diff against your
greedy. A counterexample appears in seconds if there is one.}

\subsection{Reading the Problem}
\begin{itemize}
    \item Is the answer guaranteed to exist, or must you print $-1$?
    \item How many test cases per file? Are the sums of $n$ bounded?
    \item Are indices 0-based or 1-based in the statement?
    \item Can $n = 0$ or $n = 1$? Can the graph be disconnected?
    \item Are values negative? Do they overflow \code{int}?
    \item Is the graph a tree, a DAG, or fully general?
\end{itemize}

\subsection{Bug Checklist}
Read this before rewriting anything from scratch.
\begin{itemize}
    \item \code{int} overflow: \code{(ll)a * b}, not \code{a * b}.
    \item Arrays not reset between test cases.
    \item Off-by-one: inclusive vs half-open ranges.
    \item 0-indexed code fed 1-indexed input.
    \item \code{n == 1}, empty input, all-equal values.
    \item Recursion depth on a path graph.
    \item \code{cin} left a newline before \code{getline}.
    \item Comparator not a strict weak ordering.
    \item \code{endl} inside a tight output loop (slow).
\end{itemize}
\warn{Clearing globals with \code{memset} across $10^5$ test cases is
itself $O(t \cdot n)$. Reset only the indices you actually touched.}

\subsection{Time Management}
\begin{itemize}
    \item Read \emph{every} problem before starting. The easiest is
          rarely problem A.
    \item Sort by how many teams have solved it, when the scoreboard
          shows that.
    \item Stuck for 20 minutes with no idea? Switch problems and come
          back.
    \item Write the brute force first when the statement is confusing ---
          it clarifies the problem and doubles as a stress-test oracle.
\end{itemize}

\end{multicols}

\vfill
\begin{center}
{\footnotesize\color{cmt}
CP Reference Book --- C++ Edition \quad\textbullet\quad
A Python edition with the same coverage ships alongside this one.\par}
\end{center}

\end{document}
